
\documentclass{article}

% NeurIPS 2026 style: single-blind workshop mode (authors visible when added).
% natbib is loaded by the style file; options must be passed ahead of it.
\PassOptionsToPackage{numbers,sort&compress}{natbib}
\usepackage[sglblindworkshop]{neurips_2026}
% \workshoptitle{...} % surfaces in the footer only once the [final] option is added

\usepackage[utf8]{inputenc}
\usepackage[T1]{fontenc}
\usepackage{microtype}

% Math & symbols
\usepackage{amsmath,amssymb}

% Graphics, tables, code
\usepackage{graphicx}
\usepackage{booktabs}
\usepackage{array}
\usepackage{listings}
\usepackage{xcolor}

% Diagrams
\usepackage{tikz}
\usetikzlibrary{positioning,fit,backgrounds,arrows.meta,calc,shadows}

% Links (loaded after the style file; cleveref must come last)
\usepackage[hidelinks]{hyperref}
\usepackage{cleveref}


% Convenience macros (mirrors main.tex)
\usepackage{xspace}
\newcommand{\rlm}{\textsc{RLM}\xspace}
\newcommand{\rlms}{\textsc{RLM}s\xspace}
\newcommand{\sh}{\textsc{Self-Harness}\xspace}
\newcommand{\sid}[1]{\textbf{\texttt{#1}}}

% ---- styling for the RLM edit-space figure (fig:surfaces) ----
\definecolor{editgreen}{HTML}{2E7D5B}
\definecolor{editfill}{HTML}{E7F5EE}
\definecolor{fixgrey}{HTML}{8A8F98}
\definecolor{fixfill}{HTML}{F1F2F4}
\definecolor{envtint}{HTML}{EAF1FB}
\definecolor{subtint}{HTML}{FCF1E7}
\definecolor{inkcol}{HTML}{2B2F36}
\newcommand{\cardhd}{\bfseries}
\newcommand{\cardsb}{\normalfont\scriptsize\color{black!62}}
\tikzset{
  penbadge/.pic={
    \fill[editgreen] (0,0) circle (0.095);
    \draw[white, line width=0.55pt, line cap=round] (-0.045,-0.045) -- (0.03,0.03);
    \fill[white] (-0.052,-0.052) circle (0.017);
  },
  lockbadge/.pic={
    \fill[fixgrey] (0,0) circle (0.095);
    \draw[white, line width=0.6pt] (-0.032,0.006) arc (180:0:0.032);
    \fill[white, rounded corners=0.2pt] (-0.05,-0.05) rectangle (0.05,0.014);
  },
}

% The initial harness: the sparse floor the optimization loop starts from.
\newcommand{\hzero}{\ensuremath{\mathrm{H}_{0}}\xspace}
% Figure_1.tex is a tall float; let it take more of a text page so it lands near
% the discussion instead of being deferred to a page of its own.
\renewcommand{\topfraction}{0.9}
\renewcommand{\textfraction}{0.07}
\renewcommand{\floatpagefraction}{0.85}

% ---- code-listing style for the initial-harness figure (Figure_1.tex) ----
\definecolor{lstbg}{HTML}{FAFAFB}
\definecolor{lstframe}{HTML}{D5D8DD}
\definecolor{lstkw}{HTML}{1052B8}
\definecolor{lststr}{HTML}{B25A05}
\definecolor{lstnum}{HTML}{9AA0A8}
\definecolor{lstcom}{HTML}{6B7280}
\lstdefinestyle{harness}{
  language=Python,
  basicstyle=\ttfamily\scriptsize\linespread{0.88}\selectfont,
  keywordstyle=\color{lstkw},
  stringstyle=\color{lststr},
  commentstyle=\color{lstcom},
  deletekeywords={None,True,False},
  emph={str,list,dict,int,tuple},
  emphstyle=\color{lstkw},
  numbers=left,
  numberstyle=\ttfamily\tiny\color{lstnum},
  numbersep=9pt,
  backgroundcolor=\color{lstbg},
  frame=single,
  rulecolor=\color{lstframe},
  framesep=5pt,
  breaklines=true,
  breakindent=14pt,
  breakautoindent=false,
  columns=fullflexible,
  keepspaces=true,
  showstringspaces=false,
  xleftmargin=20pt,
  xrightmargin=1pt,
  aboveskip=2pt,
  belowskip=2pt,
}

\title{Self-Harnessing Recursive Language Models}
% No authors for now; sglblindworkshop mode renders them when an \author block is added.
% \author{William Stanford\\ \small\texttt{william.stanford.c@gmail.com}}

\begin{document}

\maketitle

\begin{abstract}
Recursive language models (\rlms) let a bounded-context model operate over inputs larger
than its context window by offloading the input, decomposing it, invoking recursive model
calls, and aggregating their outputs~\citep{zhang2025rlm}. Vanilla RLMs suffer from many failure
cases, necessitating the fine-tuning LLMs and/or engineering harnesses by hand to overcome 
their limitations. We propose adapting the \sh{} framework of \citet{zhang2026selfharness} to \rlms, so that a fixed model improves its own recursive harness from evidence produced while running: it mines recurring failure mechanisms from verifier-grounded execution traces, proposes minimal edits to declared harness surfaces, and promotes only edits that survive non-regressive
validation. The optimized harness is frozen and evaluated on instances $8$--$32\times$
longer than any seen during optimization and in new target environments. Results forthcoming.
\end{abstract}

\section{Introduction}
Recursive language models (\rlms) reframe inference as a program in which a language
model inspects its own context, decomposes it, and invokes fresh copies of itself on the
pieces~\citep{zhang2025rlm}. Every call stays inside the native context window, letting a
bounded-context model reason over prompts one to two orders of magnitude larger than that
window. The mechanism behind these gains is distributional rather than architectural:
the harness shapes each call so that a task which is out-of-distribution as a whole is
handled through calls that are individually in-distribution~\citep{zhang2026harness}.

Performance therefore depends not only on the model but on the harness that determines
when and how recursion occurs. However, the RLM harness grants the capacity to recurse without
guaranteeing good recursion. Documented failures are recurring and model-specific:
policies collapse the entire problem into a single sub-call~\citep{zhang2026harness},
deeper recursion degrades accuracy while inflating cost~\citep{wang2026overthink}, and
recursion halts by heuristic rather than by evidence that further work would not
help~\citep{guha2026termination}. Yet today, models are fine-tuned to function within
the RLM harness more effectively~\citep{zhang2026harness,yang2026recursive}, or the RLM
harness is engineered by hand~\citep{alizadeh2026srlm,roy2026ycombinator,lumer2026rah}.
An alternative is to let the harness optimize itself: the model inspects its own
execution trajectories and proposes edits to the scaffold that produced
them~\citep{weng2026harness,lee2026rhi,zhang2026selfharness}. Because the evidence
driving each edit comes from the deployed model on the deployed task distribution, the
resulting harness can adapt both to the structure of the task and to model-specific
capabilities, a fit that no single hand-designed harness appears to provide
universally~\citep{zhang2026selfharness,lee2026rhi,gupta2026nosuperior}.

We ask whether a fixed language model can instead improve its own RLM harness from
evidence produced during execution:
\begin{quote}
\emph{Can a language model identify recurring failures in its own recursive behavior,
convert them into targeted harness edits, and discover an orchestration harness that
transfers across input lengths and task environments?}
\end{quote}

\section{Contributions}

Our contributions are:

\begin{enumerate}
\item \textbf{\sh{} for recursive language models.}
We adapt \sh{} to \rlms by declaring nine editable harness surfaces keyed to phases of the \rlm{} turn loop: the REPL contract, the decomposition, execution, verification and recovery instructions, the runtime policy, the metadata function, the REPL helpers, and the answer middleware (\cref{fig:edit-surfaces}). A fixed model mines recurring failure mechanisms from structured summaries of its recursive traces, proposes targeted edits, and validates them without weight updates or a stronger external model. Optimization starts from a deliberately sparse initial harness in which seven of the nine surfaces are empty, disabled, or a single generic line.

\item \textbf{Evidence of length generalization.}
We optimize only on short source instances, freeze the harness, and evaluate it on source instances $8$--$32\times$ longer, testing whether self-discovered harness improvements transfer across input length.

\item \textbf{Evidence of cross-environment transfer.}
We apply the same frozen harness to the short and long splits of an unseen target environment. This tests whether the recovered strategy transfers beyond the environment that produced it and whether it survives simultaneous changes in task environment and input length.
\end{enumerate}

\section{Methods}

\subsection{Model and Harness}
\label{sec:model-and-harness}
We use a single frozen language model, with the same model and decoding configuration serving the root and all recursive sub-calls. Three invariants of the initial reference \rlm{} implementation of \citet{zhang2025rlm}, remain fixed. 
\begin{enumerate}
\item \textbf{Prompt-as-variable.} The prompt lives in the environment as a REPL variable and is never copied into the root context.
\item \textbf{Programmatic sub-calls.} Sub-calls are issued by code, in loops, over slices of the prompt.
\item \textbf{Outputs-in-variables.} The final answer is accumulated in REPL variable and returned from it.
\end{enumerate}
Their runtime is extended with four injection points: an injected metadata function, answer-protocol middleware, sub-call retry and validation, and enforceable batch caps. Each defaults to the shipped behavior, so an unconfigured harness is byte-identical to the reference implementation. These are the seams through which the declared surfaces act; model weights, external tools, and the evaluator remain fixed.

Following \citet{zhang2026selfharness}, whose initial harness declares nine builder functions as configuration points of the agent loop, we declare \textbf{nine editable surfaces keyed to phases of the \rlm{} turn loop}, each a builder function in a single harness-definition module. \Cref{fig:edit-surfaces} shows the initial harness \hzero{} in full: each surface with the value it returns at the starting point of optimization (the condition reported as B1 in \cref{sec:baselines}). 

Declaring the surfaces from the loop's phase structure, rather than selecting them from failures already documented for \rlms{}, is deliberate. A surface set chosen by reading a published catalogue of failures would make it impossible at analysis time to separate what the optimization loop discovered from what its designer already knew. Surface \emph{selection} is part of the fixed scaffold shared by every condition, exactly as it is in \citet{zhang2026selfharness}: the model proposes bounded edits to declared surfaces and does not choose which surfaces exist.

\hzero{} is deliberately sparse. Seven of its nine surfaces return nothing, a disabled policy, or a single generic line, so the orchestration guidance a capable harness eventually carries is guidance the loop has to recover from its own traces rather than reread from its own system prompt.

% Figure 1: the initial harness H0 and its nine editable surfaces.
% Body only; the `harness' listing style, its colours, and \hzero live in the
% preamble of proposal.tex, matching how Figure_2.tex relies on the tikz styling
% defined there.
\begin{figure}[!tbp]
\centering
\begin{lstlisting}[style=harness]
# shrlm/rlm_harness.py: the nine surfaces Self-Harness may edit (S1 to S9).

MINIMAL_REPL_CONTRACT = """
You are a Recursive Language Model (RLM): a language model whose prompt-related
information lives in a Python REPL that you drive turn by turn.
Available in the REPL:
- `context`: the information related to the prompt (`str` or `list[str]`).
- `llm_query(prompt, model=None) -> str`; `llm_query_batched(prompts, ...)`.
- `rlm_query(...)` / `rlm_query_batched(...)`: recursive RLM sub-calls.
- `SHOW_VARS() -> str`: list every variable currently in the REPL.
- `answer`: dict {"content": "", "ready": False}; set both to submit.
{custom_tools_section}
Write code in ```repl``` blocks, one per turn; the REPL persists. Only
`print(...)` output is shown back. Outputs over ~20K characters are truncated.
"""

def build_repl_contract() -> str:                                    # S1
    return MINIMAL_REPL_CONTRACT

def build_decomposition_instruction() -> str:                        # S2
    return ("Start by probing `context` to understand what you have, "
            "then decide how the task breaks into steps.")

def build_execution_instruction() -> str:                            # S3
    return ("On each turn, run one block of code, print what you need "
            "to see, and build the result up in REPL variables.")

def build_verification_instruction() -> str:                         # S4
    return 'Check your answer before setting `answer["ready"] = True`.'

def build_recovery_instruction() -> str:                             # S5
    return ("If a sub-call errors or returns something you cannot use, "
            "adjust the approach and try again.")

def build_runtime_policy() -> dict[str, Any]:                        # S6
    return {
        "enabled": False,
        "max_prompt_chars": None,
        "max_batch_width": None,
        "max_depth": None,
        "retry_on_syntax_error": None,
        "max_retries": None,
        "validate_sub_output": None,
    }

def build_metadata(                                                  # S7
    stdout: str,
    repl_inventory: dict[str, tuple[str, int]],
    max_character_length: int = DEFAULT_MAX_CHARACTER_LENGTH,
) -> str:
    return default_metadata_builder(stdout, repl_inventory,
                                    max_character_length)

def build_repl_helpers() -> dict[str, Any]:                          # S8
    return {}

def build_sub_repl_helpers() -> dict[str, Any]:                      # S8
    return {}

def build_answer_middleware() -> AnswerMiddlewareFn:                 # S9
    return accept_answer
\end{lstlisting}
\caption{Initial harness \hzero{} and the editable interface used as the starting
point for \sh{}: nine declared surfaces, one builder function each, keyed to phases
of the \rlm{} turn loop. The harness is intentionally kept sparse, with seven of the
nine returning nothing, a disabled policy, or a single generic line, so that the
orchestration strategy a tuned harness carries is something the optimization loop
must recover from its own execution traces. \sid{S6}, \sid{S7} and \sid{S9} required
a runtime injection point in the reference implementation of \citet{zhang2025rlm};
the rest are reached through its existing prompt and tool configuration. A proposal
modifies exactly one surface.}
\label{fig:edit-surfaces}
\end{figure}



\subsection{Environments and splits.}
We evaluate on two environments from the length-generalization suite of \citet{zhang2026harness}.

\begin{itemize}
\item \textbf{GraphWalks}~\citep{openai2025graphwalks}, multi-hop reachability and traversal queries over a graph presented in context. A sub-call asking for reachability within an induced subgraph has a checkable answer, so a child can be verified directly.
\item \textbf{OOLONG-Pairs}, a modified variant of OOLONG~\citep{bertsch2025oolong} that asks for pairs of elements satisfying a given constraint. A candidate pair either satisfies the constraint or does not, which is checkable in isolation.
\end{itemize}

\noindent Both require decomposing an input into independently solvable units and aggregating their outputs, but they differ in surface task and data distribution, which is the preregistered structural criterion by which they were chosen. Each provides matched short and long instances, with long inputs approximately $8$--$32\times$ larger, and a deterministic verifier that we use unmodified. One is preregistered as the source environment and the other serves as the target.

Both environments were selected because they additionally admit \emph{synthesized sub-verifiers}: the sub-problems a decomposition produces are themselves checkable without reference to the root answer. This is what makes the failure attribution of \cref{sec:sho} more than a model judgment. When the root answer is wrong, running the sub-verifier on each child separates a root failure, in which correct children were aggregated into a wrong answer, from a child failure, in which a sub-call returned a wrong local result that the root faithfully combined. The two implicate different editable surfaces, and without sub-verification the distinction rests on the proposer's reading of a trace rather than on a checkable outcome.

Only the source environment is used for optimization. Its short split is partitioned into held-in, held-out, and test sets. Held-in traces support weakness mining, the held-out split serves only as the promotion gate, and the source-short test and source-long test remain inaccessible until the harness is frozen. No target-environment instances, traces, verifier outcomes, or aggregate results are exposed during optimization. The frozen harness is subsequently evaluated on both the short and long target splits.

\input{Figure_2.tex}

\subsection{Self-Harness Optimization}
\label{sec:sho}

% \SC{Clarification: the process still has a very small meta-harness that informs the current harness of the task for which it is being summoned (solve a dataset instance or improve itself), correct?}

% \SC{A diagram similar to the one from the self-harness paper that highlighs the  RLM specificities would be useful here and reusable for the paper.The more detailed the better.}

% \SC{It seems to me that this could be more detailed (or perhaps simply clearer / more structured) when listing the components that go in the pipeline. This would be useful to spread the work.}

We adapt the \sh{} framework of \citet{zhang2026selfharness} to \rlms. The model weights remain fixed throughout, and the same model is used both to execute tasks and to propose harness changes. The editable surfaces are the nine declared in \cref{fig:edit-surfaces}, and a proposal modifies exactly one of them. The model may not modify the evaluator, external tools, or any of the three invariants of \cref{sec:model-and-harness}. One further surface is permanently excluded: allowing the proposer to truncate or pre-summarize the prompt variable into root context. It would likely show held-in gains, and it violates the first invariant by turning the \rlm{} back into a compaction agent. Each optimization round has three stages (\cref{fig:sh-loop}): \textbf{Weakness Mining}, which converts failed runs into a verifier-grounded evidence bundle of recurring failure patterns; \textbf{Harness Proposal}, which turns mined patterns into a small set of minimal candidate edits; and \textbf{Proposal Validation}, which promotes only edits that survive regression testing on data never shown to the proposer.

\subsubsection{Weakness Mining}
\label{sec:weakness-mining}

Weakness mining runs the current harness on short held-in instances from a source environment, recording the verifier outcome and full recursive execution trace of every run, including runs terminated by resource limits. Each recorded sub-call is then scored post hoc by the environment's synthesized sub-verifier; because sub-problems are authored by the model itself, a sub-call may be correct, incorrect, or uncheckable, and the level at which the error signal first appears is treated as a checked fact only when at least one sub-call was checkable. Records with no checkable sub-call are marked ungrounded and their failing level falls to the model's judgment. Mining runs in both conditions — with the sub-verifier and with it withheld, in which case every record is ungrounded — and because sub-verification is applied post hoc to recorded traces, both conditions can mine the same runs, holding execution fixed while only the grounding signal varies; the comparison is the ablation of \cref{app:ablations}. Each failure is converted into a structured record: the verifier-level cause and the failing level are derived from the verifier and sub-verifier, while the causal status of the associated agent behavior and the implicated mechanism are judged by the same fixed model from a bounded, mechanically compressed digest of the trace. These judgments are validated against a closed vocabulary of mechanisms, each mapping to exactly one editable surface, and a judgment that never validates is recorded as unattributed rather than coerced. Failures are then clustered by exact agreement on the verifier-grounded signature
\begin{equation*}
\phi(r_i)=(\text{verifier cause},\ \text{failing level},\ \text{causal status},\ \text{agent mechanism}),
\end{equation*}
and the resulting evidence bundle ranks each recurring pattern by the number of distinct instances it affects and its estimated actionability, carrying cluster size, mechanically computed shared symptoms, and representative trace excerpts. The bundle describes weaknesses without prescribing edits, and all artifacts behind it, the traces, per-child labels, digests, and judgments, are persisted for audit.

\subsubsection{Harness Proposal}
\label{sec:harness-proposal}

Harness proposal provides the same fixed model with the verifier- and sub-verifier-grounded failure patterns of the evidence bundle, a record of passing behaviors that should be preserved, and a summary of previously attempted edits, and asks it to generate several distinct, minimal candidate edits, excluding patterns not plausibly addressable by an editable surface. Each proposal must target one mined failure pattern, modify one declared harness surface, state its predicted behavioral effect, and identify possible regressions.

\subsubsection{Proposal Validation}
\label{sec:proposal-validation}

Proposal validation evaluates the incumbent and each candidate on both the source held-in split and a disjoint short held-out split that is never shown to the proposer, with a preregistered number of repetitions per instance; pass counts are aggregated across repeats. Candidates that fail structural checks --- a stale base harness, zero or multiple surfaces changed, a harness-invariant violation --- are rejected without evaluation. A candidate is accepted only if $\Delta_{\mathrm{in}} \ge -\tau_{\mathrm{reg}}$, $\Delta_{\mathrm{ho}} \ge -\tau_{\mathrm{reg}}$, and $\max(\Delta_{\mathrm{in}}, \Delta_{\mathrm{ho}}) > \tau_{\mathrm{imp}}$, where the $\Delta$ are aggregated pass-count differences against the incumbent on each split and the tolerances are preregistered from pilot variance (zero recovers the strict rule of \citet{zhang2026selfharness}), and only if its mean sub-call count and cost fall within a preregistered band relative to the incumbent. Because the runtime policy is itself an editable surface, evaluation limits (recursion depth, iterations, per-run and per-candidate budgets) are enforced by the experiment outside the harness; a candidate's policy may tighten them but never exceed them. Unlike \citet{zhang2026selfharness}, when multiple accepted edits touch disjoint surfaces, we re-evaluate the merged harness under the same rule before promotion; if the merge fails, the round promotes nothing, since falling back to individually accepted edits after observing the merge result would be post-hoc selection outside the preregistered rule. Optimization stops after a fixed number of rounds or several consecutive rounds without a promotion.

\subsection{Baselines and Evaluation}
\label{sec:baselines}
We evaluate the Self-Harnessed RLM (SH-RLM) against two baselines.

\begin{itemize}
\item \textbf{B1: initial \rlm{} (\hzero{}).} The unmodified reference harness with every surface of \cref{fig:edit-surfaces} at its \hzero{} default, and the starting point of optimization.
\item \textbf{H1: $\lambda$-RLM~\citep{roy2026ycombinator}.} A hand-designed extension, carried over unmodified, that replaces free-form recursive code generation with a typed functional runtime and invokes the model only on bounded leaf sub-problems. It measures \sh{} against human harness engineering rather than against an unmodified starting point.
\end{itemize}

\noindent One further condition relaxes the fixed-weight constraint and is reported separately.

\begin{itemize}
\item \textbf{F1: fine-tuned \rlm{}.} The same backbone RL-trained inside the initial harness on the source short split, following \citet{zhang2026harness} (\cref{app:finetune}). It bounds how much of the residual gap after harness optimization is attributable to model capability rather than orchestration. Because weight training and harness optimization incur non-comparable costs, F1 is a reference point rather than a matched comparison, and is excluded from the primary comparisons.
\end{itemize}

Each condition is evaluated on four untouched categories of test set, with the two target categories reported separately for each target environment:
\begin{enumerate}
\item \textbf{source-short}, measuring improvement at the optimization length;
\item \textbf{source-long}, measuring length generalization;
\item \textbf{target-short}, measuring cross-environment transfer without a length shift; and
\item \textbf{target-long}, measuring cross-environment transfer under an additional length shift.
\end{enumerate}

The primary comparisons are SH-RLM versus B1 on source-long, target-short, and target-long, with SH-RLM versus H1 on the same three sets indicating whether a self-discovered harness is competitive with a hand-designed one. Improvement on source-long indicates that the learned harness changes survive a length shift within the optimization environment. Improvement on target-short provides the cleanest evidence that the edits capture a transferable compositional strategy rather than a source-specific rule. Improvement on target-long tests whether that transfer remains effective when environment and input length change simultaneously.

\subsection{Analysis}
\label{sec:analysis}

The primary metric is verifier accuracy, reported separately for all four test
sets. Results are averaged over repeated seeded runs and accompanied by
bootstrap confidence intervals over task instances. B1 and SH-RLM are compared
using a paired per-instance test; the exact test and repetition-to-instance
aggregation rule will be preregistered.

Secondary efficiency metrics are total input and output tokens, recursive-call
count, maximum recursion depth, and accuracy per million tokens. Trace analysis
tests whether optimization changes the mechanisms it was intended to repair. In
particular, we report the frequency of each mined failure pattern before and
after optimization and measure whole-input sub-call collapse, defined as a run
that delegates most of the input to one child or performs no meaningful
decomposition. Using the sub-verifiers, we also report the share of failures
attributable to the root and to the children in each condition and split,
which shows whether optimization moved errors down the call tree, repaired
them, or merely relocated them.

Two ablations, one withholding the sub-verifier signal from the optimization loop
and one removing each promoted edit from the final harness, are specified in
\cref{app:ablations}. Additionally, we include an optional analysis on 
alignment drift in \cref{app:alignment}.

\section{Feasibility and Cost}
\label{sec:cost}

The \sh{} optimization loop runs only on the source environment. Let $n_{\mathrm{in}}$ and $n_{\mathrm{ho}}$ denote the source held-in and held-out sizes, $m$ the number of mining repetitions, $v$ the validation repetitions, and $K$ the proposal width. Excluding occasional merged-candidate checks, each optimization round requires
\begin{equation*}
R_{\mathrm{round}}
=
m\,n_{\mathrm{in}}
+
v(K+1)(n_{\mathrm{in}}+n_{\mathrm{ho}})
\end{equation*}
runs. The first term collects source held-in traces for weakness mining; the second evaluates the current harness and $K$ candidate harnesses on the two source optimization splits. Optimization is capped at $T$ rounds and stops early after a preregistered number of rounds without promotion.

At $n_{\mathrm{in}}=24$, $n_{\mathrm{ho}}=40$, $m=2$, $v=4$, and $K=4$, this is $1{,}328$ runs per round and $19{,}920$ short runs over $T=15$ rounds (\sh{} used 15, 18, and 21, rounds for the three model families tested). Taking three passes over the stored context plus root overhead gives roughly $1.2\times10^{5}$ tokens for a typical short run, so optimization costs approximately $2.4\times10^{9}$ tokens.

After optimization, the three fixed-weight conditions B1, H1, and SH-RLM are evaluated on the source-short, source-long, target-short, and target-long test sets; F1 is budgeted separately (\cref{app:finetune}). The target environment therefore adds final-evaluation cost but no mining, proposal, or candidate-validation cost. Total cost can be expressed as
\begin{equation*}
C_{\mathrm{total}}
=
C_{\mathrm{source\ optimization}}
+
C_{\mathrm{source\ evaluation}}
+
C_{\mathrm{target\ evaluation}}.
\end{equation*}

Evaluating three conditions on two environments at $40$ short-test and $150$ long-test instances with $3$ repetitions is $720$ short and $2{,}700$ long runs. The long runs dominate: at inputs $8$--$32\times$ larger, they average roughly $1.2\times10^{6}$ tokens each, giving approximately $3.2\times10^{9}$ tokens for final evaluation against $9\times10^{7}$ for the short tests. The project total is therefore approximately $5$--$6\times10^{9}$ tokens. Served locally that is roughly $650$ H100-hours; purchased as hosted inference at current open-weights rates it is approximately \$$1{,}200$--\$$2{,}000$. The sub-verification ablation of \cref{app:ablations} adds one further optimization run, or approximately $2.4\times10^{9}$ tokens, and is budgeted as contingent.

Cost is dominated by long-instance evaluation and by repeated reads of the offloaded context during recursion, and the passes-over-context multiplier is the dominant uncertainty: a factor-of-two error moves the total to between $3\times10^{9}$ and $1.1\times10^{10}$ tokens. We will run a small pilot in both environments before optimization to measure tokens per run, recursive calls, and effective passes over the stored context, and these measurements will fix the final test sizes and token budget before the first harness edit is proposed. If needed, we can reduce long-test sample size in all four evaluation conditions to reduce cost.

  \section{Timeline}

\paragraph{Completed work.} Three components of the system are already built and working. Both evaluation environments are implemented with their deterministic verifiers and synthesized sub-verifiers integrated. The initial \rlm{} is implemented with the nine editable harness surfaces of \cref{fig:edit-surfaces} declared and their runtime injection points in place. The first stage of the \sh{} loop, \textbf{Weakness Mining}, is operational: recursive execution traces are collected, converted into structured failure records, and clustered by verifier-grounded signature.

The remaining work is designed for completion over six weeks. Each week is split into three chunks that are independent of one another and can proceed in parallel; dependencies run only between weeks.

\begin{description}

\item[Week 6: Preregistration and baselines; begin the \sh{} loop; write introduction.]
(i)~Preregister the source and target environment assignment, task-structure criterion, model checkpoint, split sizes, evaluation repetitions, promotion rule, and inference budgets, and reproduce baseline behavior on a small sample. (ii) Implement \textbf{Harness Proposal}: failure pattern selection, summary of previous edits, a history of behaviors to preserve, parallel harness edit proposal. (iii)~Finish the introduction.

%\SC{This is already started but matches with week 6 in Algoverse's numeration (the 02/08/2026 is week 5)}

\item[Week 7: Pilot and calibration; continue the \sh{} loop; write related work.]
(i)~Run pilot instances from both environments to estimate token cost, latency, context rereading, and baseline variance; define the promotion gate's empirical noise margin from repeated source-short evaluations; finalize test sizes and apply any preregistered cost de-scoping. (ii)~Implement \textbf{Proposal Validation}: eval on held in/out sets for regression testing, promotion rule based on cost and performance, merged evaluation. (iii)~Write the related work.

\item[Week 8: Validate the loop; build evaluation tooling; write methods.]
(i)~Test the complete \sh{} loop on a small development subset excluded from the final experiment, and verify that runs are reproducible from saved configurations. (ii)~Integrate the hand-designed $\lambda$-RLM baseline (H1) and build the evaluation and analysis tooling: matched-budget configurations, verifier-accuracy aggregation, confidence intervals, paired tests, and efficiency metrics. (iii)~Write the methods section.

\item[Week 9: Source-environment optimization and freeze; write ablations and discussion.]
(i)~Run \sh{} on the source held-in and held-out splits, inspecting logs for infrastructure failures without manually changing proposed edits or promotion decisions, and freeze the final harness once the maximum-round or patience criterion is reached; no test split is opened before the freeze. (ii)~Dry-run the four-way evaluation pipeline on the optimization splits so evaluation starts immediately at freeze, and run the optional alignment-drift probe across the accumulating harness lineage. (iii)~Draft the ablations appendix and begin the discussion from optimization-split trace case studies.

\item[Week 10: Fixed-weight evaluation; analysis; write results.]
(i)~Run the four-way evaluation of B1, H1, and SH-RLM on the untouched test sets under matched inference budgets, with leave-one-edit-out ablations where budget permits and a published fine-tuned checkpoint if one is available. (ii)~As results arrive, compute verifier accuracy, confidence intervals, paired comparisons, token efficiency, recursion depth, call counts, and whole-input sub-call collapse. (iii)~Write the results section.

\item[Week 11: Consolidation and release; write limitations and conclusion.]
(i)~Complete any remaining long-split evaluation runs and ablations. (ii)~Compare the frequency of mined failure mechanisms before and after optimization and finish the trace case studies. (iii)~Assemble the manuscript, write limitations and conclusion, and release the split identifiers, configurations, harness lineage, proposal and validation logs, and evaluation scripts.

\end{description}

The critical path runs through chunk (i) of each week: preregistration, pilot, loop validation, harness freeze, then evaluation on data that remains untouched during optimization. The remaining chunks are parallel engineering and writing tracks that keep the freeze date off the writing's critical path.


\clearpage
\bibliographystyle{unsrtnat}
\bibliography{references}

\clearpage
\appendix

\section{Fine-tuned \rlm{} baseline (F1)}
\label{app:finetune}

F1 reproduces the weight-training arm of \citet{zhang2026harness} so that harness
optimization and weight optimization are compared on the same backbone, harness, and
environment. The recipe below follows their reported setup; any deviation forced by our
compute allocation will be recorded before training begins.

\begin{description}
  \item[Backbone.] Qwen3-30B-A3B-Instruct-2507, the model \citet{zhang2026harness}
    RL-train inside an \rlm{} harness, and the same backbone used for B1, H1, and SH-RLM so that F1
    differs from B1 only in its weights.
  \item[Algorithm.] RL with \texttt{prime-rl}: decoupled PPO with GRPO-style advantages
    and a KL penalty against the initial policy.
  \item[Reward.] The environment's deterministic verifier on the final answer, identical
    to the evaluation metric and to the promotion signal used by \sh{}, so the two
    optimizers are driven by the same outcome measure.
  \item[Rollouts.] Batch size $64$ with $4$ rollouts per sample.
  \item[Steps.] $150$ optimization steps for the length-generalization setting, with
    evaluation every $10$ steps; \citet{zhang2026harness} run $500$ steps, evaluating
    every $20$, for their cross-domain strategy experiments.
  \item[Training inputs.] Short instances only, at $8$k--$64$k tokens, drawn from the
    same source-environment short split used for harness optimization. Long instances are
    never trained on.
  \item[Evaluation.] The frozen checkpoint is evaluated on the same four untouched test
    sets as B1, H1, and SH-RLM, at $256$k--$2$M tokens on the long splits.
  \item[Hardware.] $8\times$H100 nodes.
\end{description}

Two conditions govern whether F1 is reported at all. If the published checkpoint is
obtainable, we evaluate it directly rather than retraining, and say so. If neither the
checkpoint nor the compute for training is available, F1 is dropped and its absence is
reported alongside the published figures for the same environment and backbone. F1 is
budgeted separately from the cost estimate of \cref{sec:cost}, which covers the
fixed-weight conditions only.

\section{Ablations}
\label{app:ablations}

\paragraph{Sub-verification.} Sub-verification enters the loop only through
weakness mining, so we can withhold it without changing anything else. We repeat
the full optimization run with the sub-verifier signal removed, giving the
proposer root outcomes and traces alone, and compare the resulting frozen harness
against SH-RLM on all four test sets. The comparison asks whether checkable
child-level evidence is what makes mined failure attributions actionable, or
whether the proposer recovers the same edits from traces without it. Because
sub-verification is post hoc, the same recorded runs can additionally be mined
under both conditions within a round, yielding paired evidence bundles that
differ only in the grounding signal; this trace-level comparison costs no extra
execution and complements the loop-level ablation. Because the full repeat
doubles optimization cost, it is a preregistered secondary run on the source
environment, contingent on the pilot's cost estimate holding.

\paragraph{Leave-one-edit-out.} Each promoted edit is removed from the final
harness individually and the affected test conditions are re-evaluated. These
ablations determine which edits are responsible for the final gains and whether
improvements arise from the accumulated harness rather than unrelated sampling
variation.

\section{Assessing Alignment Drift (Optional)}
\label{app:alignment}
  
Given prior work on misevolution in self-optimizing agents~\citep{shao2025misevolution},
and in particular the finding that capability-driven workflow optimization can sharply
reduce refusal rates on harmful requests despite passing task-level
validation~\citep{shao2025misevolution}, an optional extension is to evaluate SH-RLM for
alignment drift as a side effect of harness promotion. Because our promotion gate, like
that of \citet{zhang2026selfharness}, is defined only over verifier accuracy and cost,
edits that are capability-positive but compliance-increasing (e.g., ``always commit to an
answer,'' suppressed hedging, or reasoning-mode changes, which are known to shift safety
behavior in the Qwen3 family~\citep{selfjailbreak2025}) would be promoted without
detection. The extension freezes a small safety probe of a few hundred harmful and
borderline prompts~\citep{wildjailbreak2024}, runs it through the \rlm{} under B1 and
under each promoted harness in the lineage, and scores refusal and harmful-answer rates
with an automatic judge, yielding a safety trajectory plotted alongside the accuracy
trajectory. At roughly $300$ prompts across $\sim$$8$ harness states, with probe runs
far shorter than benchmark instances, this adds on the order of $10^{7}$--$10^{8}$
tokens, well under $2\%$ of the budget in \cref{sec:cost}. It would enable us to assess
whether the propose--validate--accept structure incidentally preserves alignment, as its
bounded edit surfaces might suggest, or whether the harness pathway exhibits the same
drift documented for workflow evolution; either outcome is informative, and to our
knowledge neither has been measured for \sh{}-style harness optimization. Because no
promotion decision depends on the probe, it is purely observational and does not alter
the preregistered optimization loop.

\end{document}




  